% Part: model-theory
% Chapter: interpolation
% Section: introduction

\documentclass[../../../include/open-logic-section]{subfiles}

\begin{document}

\olfileid{mod}{int}{int}

\olsection{పరిచయం}

అంతర్వేశ సిద్ధాంతం కింది ఫలితం: $\Entails !A \lif !B$ అనుకుందాం.
అప్పుడు $\Entails !A \lif !C$, $\Entails !C \lif !B$ అయ్యే
\tetoken{ఒక వాక్యం}{!!a{sentence}}~$!C$ ఉంటుంది. అంతేకాక, $!C$లోని
ప్రతి \tetoken{స్థిరసంకేతం}{!!{constant}},
\tetoken{ప్రమేయ సంకేతం}{!!{function}},
\tetoken{విధేయ సంకేతం}{!!{predicate}} ($\eq$ తప్ప) $!A$, $!B$
రెండింటిలోనూ వస్తుంది. ఆ \tetoken{వాక్యం}{!!{sentence}}~$!C$ను
$!A$, $!B$ల \emph{అంతర్వేశకం (ఇంటర్‌పోలంట్)} అంటారు.

అంతర్వేశ సిద్ధాంతం స్వతహాగా ఆసక్తికరమైన ఫలితమే. అయితే ఒక సిద్ధాంతంలో
నిర్వచనీయత గురించిన ఫలితాలను, అలాగే రెండు అవైరుధ్య సిద్ధాంతాలను
కలిపినప్పుడు అవైరుధ్య సిద్ధాంతమే ఏర్పడే షరతులను నిరూపించడానికి దాన్ని
ఉపయోగించవచ్చనేదే దాని ప్రధాన ప్రాముఖ్యం. మొదటి ఫలితాన్ని బెత్
నిర్వచనీయతా సిద్ధాంతం అని, రెండవదాన్ని రాబిన్సన్ ఉమ్మడి అవైరుధ్య
సిద్ధాంతం అని అంటారు.

\end{document}
